\section{Admissibility and meta-theory}\label{sec:admissibility-meta-theory}


This section states when a claim about a description counts. In
Definitions~\ref{prop:honest-bookkeeping}--\ref{prop:empirical-bridge} and
Propositions~\ref{prop:typed-non-collapse} and~\ref{prop:nonclaim-closure}, a
claim is a claim annotation of some \(D\in\FATCD{}\), and ``inadmissible''
means that condition (a) or (b) of Definition~\ref{prop:honest-bookkeeping}
fails, or that a record required by Definition~\ref{prop:level-profile},
\ref{prop:forgetting-lifts}, \ref{prop:visibility} or~\ref{prop:empirical-bridge}
is missing; in particular a claim other than an offer that no honest audit annotation of
\(D\) certifies is inadmissible; an offer satisfies condition (a) of Definition~\ref{prop:honest-bookkeeping} when an honest audit
records the value of \(W_i(C)\), even when that value is false. The central
rule is honest bookkeeping (Definition~\ref{prop:honest-bookkeeping}):
every claim, and every change of level, translation, visibility
assumption or empirical bridge it relies on, must be explicitly
recorded. Together with the role meanings and boundary distinctions of
Section~\ref{sec:primitive-role-architecture}, that rule governs the
claims developed below.
\begin{itemize}[topsep=2pt,itemsep=1pt]
  \item The six roles are kept apart as types (typed non-collapse); the
    paper does not claim that they are independent generators.
  \item A claim moves between levels only through recorded forgetting,
    which loses data, and recorded selected lifts, which add it.
  \item A role counts as active only with the required threshold,
    witness, level, collapse, nonclaim and audit data.
  \item An instrument decides what is visible and what is suppressed.
  \item An empirical claim enters only through a recorded bridge.
\end{itemize}
The section uses the domain \FATCD{}, the audit part \(\Audit{D}\) and
threshold attachment, which are defined in
Section~\ref{sec:fatcd-domain}; a reader checking proofs can follow the reading order at the end of
Section~\ref{subsec:organization}, keeping Definition~\ref{def:audited} at
hand while reading Definition~\ref{prop:honest-bookkeeping}. Several claim-level rules
below use records and annotations defined in
Section~\ref{sec:fatcd-domain}. The main correspondences are:
\begin{itemize}[topsep=2pt,itemsep=1pt]
  \item Definition~\ref{prop:honest-bookkeeping} and
    Definition~\ref{def:audited}: Definition~\ref{def:audited} requires every
    claim of a description to be referred to by an honest audit;
    Definition~\ref{prop:honest-bookkeeping} further requires that audit to
    certify the claim's recorded statement (for an offer, the value of
    \(W_i(C)\));
  \item Remark~\ref{prop:activation-thresholds} and
    Definition~\ref{def:threshold-attachment}: when a role reading of a
    cluster is licensed;
  \item Definition~\ref{prop:forgetting-lifts} and the lift and forgetting
    annotations of Definition~\ref{def:annotations}: how a claim changes
    level;
  \item Definition~\ref{prop:visibility} and the visibility annotations:
    what an instrument shows and suppresses;
  \item Definition~\ref{prop:empirical-bridge} and the empirical-bridge
    annotations: how an empirical claim enters.
\end{itemize} The rules keep the global
nonclaims of Remark~\ref{rem:atlas-global-nonclaims} in force. Of this section, only
Definitions~\ref{prop:honest-bookkeeping}, \ref{prop:level-profile},
\ref{prop:forgetting-lifts}, \ref{prop:visibility}
and~\ref{prop:empirical-bridge} and Proposition~\ref{prop:nonclaim-closure}
are cited in Sections~\ref{sec:fatcd-domain}--\ref{sec:scoped-exact-six}, the
last only in Section~\ref{sec:integrated-stress}. The other results of this section
fix how claims are stated; several return in
Sections~\ref{sec:discussion}--\ref{sec:future-work}.

\subsection{Honest bookkeeping as an admissibility regime}\label{subsec:honest-bookkeeping}

The audit condition of Definition~\ref{def:audited} records enough
data for each claim, promotion, translation, lift, forgetting,
visibility, and empirical-bridge assertion to be interrogated on its
own terms. The following definition stipulates that these record
requirements jointly constitute the admissibility regime, in a spirit
analogous to the axiomatic specification discipline of Hoare
logic~\citep{Hoare1969}.

\begin{definition}[Admissible claim under honest bookkeeping]\label{prop:honest-bookkeeping}
Honest bookkeeping is the admissibility regime for the claim annotations
carried by a description \(D \in \FATCD{}\). A claim annotation of \(D\) is
\emph{admissible} if and only if (a) some honest audit annotation of
\(D\) refers to it and certifies the formula recorded as its statement (for an offer, \(W_i(C)\) or
\(\neg W_i(C)\) according to the recorded result, Definition~\ref{def:audited}),
and (b) every assertion below has an explicit record carried in \(\Audit{D}\)
in the sense of Definition~\ref{def:annotations}:
\begin{itemize}[topsep=2pt,itemsep=1pt]
  \item the claim itself, with its statement, hypotheses, and scope;
  \item every promotion, with its source level, target level, added data,
    losses, and preserved nonclaims;
  \item every translation, with its source context, target context,
    preserved data, suppressed data, and claim strength;
  \item every selected lift, marked selected and non-unique unless uniqueness for its specified source data and target class is proved and recorded;
  \item every forgetting map, with its record of losses, empty only when
    losslessness is proved and recorded;
  \item every visibility claim, with both its visible and its suppressed
    data;
  \item every empirical-bridge assertion, with its substrate, observable,
    instrument, level map, threshold-witness relation, and
    suppressed-structure nonclaims;
  \item every cross-level claim, with a level annotation targeting it and
    translation annotations for each pair of levels recorded on the records it
    names, as in Definition~\ref{prop:level-profile}.
\end{itemize}
Accordingly, unrecorded claims, implicit promotions, silent translations,
asserted-but-unselected lifts, lossless forgetting without proof,
visibility overreads, and bare empirical-realization assertions are all
inadmissible.
\end{definition}

\begin{remark}[Why the regime is necessary]\label{rem:why-honest-bookkeeping}
The bookkeeping requirement is a structural check, not a stylistic
one. Without explicit records, passive observation can be misread as
\Psix{}, mere packaging as \Pone{} descent, and route failure as
\Pthree{} mismatch. The regime prevents those misreadings from
entering the theorem chain.
\end{remark}


\subsection{Typed non-collapse}\label{subsec:typed-non-collapse}

The regime does not include flat primitive independence --- the claim
that \Pone{}, \ldots, \Psix{} are independent generators of a global
algebra. Its meta-theoretic counterpart here is typed non-collapse.
Proposition~\ref{prop:boundary-distinctions} excludes each forbidden collapse
of Table~\ref{tab:boundary-matrix}; Proposition~\ref{prop:typed-non-collapse}
states the same exclusion as a rule of the admissibility regime of
Definition~\ref{prop:honest-bookkeeping} and adds that each row's permitted
interaction is not thereby excluded. Each row is checked in one direction, for the role it credits (listed in the
table after the proof of Proposition~\ref{prop:boundary-distinctions}); the
row label names the pair without order.

\begin{proposition}[Typed non-collapse]\label{prop:typed-non-collapse}
Under honest bookkeeping (Definition~\ref{prop:honest-bookkeeping}) and
the alignments of Remark~\ref{rem:repaired-alignments}, consider each
forbidden-collapse row of Table~\ref{tab:boundary-matrix}. Write \(P_c\)
for the role that row would credit and \(X_b\) for the bare partner datum
it borrows, as listed in the last column of Table~\ref{tab:boundary-matrix}
and presented as a boundary record in the sense of
Definition~\ref{def:annotations}.
For every \(D\in\FATCD{}\), no honest audit annotation of \(D\)
certifies \(W_c(C)\) for a cluster \(C\) formed only from \(X_b\); the ``permitted interaction'' of the same row is
not excluded by this proposition. More generally, for each \(k\le6\), no
honest audit annotation of \(D\) certifies \(W_k(C)\) for a cluster \(C\)
lacking any record that the \(W_k\) row requires a witness cluster to contain (among them, for \(k=1\): a map or partial-map record of declared use update, or a relation record of declared use quotient; \(k=2\): a macro-specification predicate record or a realization relation; \(k=3\): two defined path or derivation records; \(k=4\): a law-bearing stage, order, transition or refinement record; \(k=5\): a relation record of declared use quotient, a packaging map record or a check record; \(k=6\): an event record, a trace record or a replay or check record); in particular each row's exclusion holds with its two
roles exchanged whenever the exchanged datum lacks that record.
\end{proposition}

\begin{proof}
Fix a forbidden-collapse row and its credited role \(P_c\) and borrowed datum \(X_b\). The row records a ``forbidden collapse'': an
identification asserting that the content of its two roles is the same
closure role, so that one
projection may be read off the other without further data. To be
admissible under Definition~\ref{prop:honest-bookkeeping}, such an
identification would have to enter the theorem chain as a claim with a
recorded statement and scope. The role meanings are distinct, and
Proposition~\ref{prop:boundary-distinctions} checks, for each
forbidden-collapse row, a field required by the credited witness that the
named partner datum alone does not supply. An audit certifying the credited
projection solely from that datum would therefore certify a conclusion
unsupported by its recorded check and would not be honest under
Definition~\ref{def:audited}: by that definition the check-language conclusion
\(W_c(C)\) of an honest audit would be true in \(D\), whereas \(W_c(C)\) is false by
Proposition~\ref{prop:boundary-distinctions}. Independently supplied data may satisfy both
witness conditions; the proposition excludes the claimed automatic
identification, not that joint occurrence. Hence the collapse is
inadmissible. More generally, if \(C\) lacks a record required by the \(W_k\) row, that row gives \(W_k(C)=\mathrm{false}\). Definition~\ref{def:audited}(c3) therefore excludes an honest audit certifying \(W_k(C)\); this also proves the stated exchanged-role exclusions.

It remains to check that ruling out the collapse does not also rule out
the row's ``permitted interaction''. The permitted interaction does not
identify \(P_i\) with \(P_j\); it records a legitimate joint occurrence
of the two distinct roles in a single description --- for instance,
\Pone{} descent audited by \Psix{}, or \Pthree{} route terms indexed by
\Pfour{} stages, as in Remark~\ref{rem:not-global-independence}. Such
an interaction carries its own honest-bookkeeping record for each role
separately and asserts no identification. Nothing in the argument above
touches it: the inadmissibility derived for the collapse rests on the unsupported
inference from the bare partner datum to the credited witness, which the
permitted interaction never makes. Therefore the permitted interaction remains admissible
whenever its own bookkeeping conditions are met.
\end{proof}

\begin{remark}[Typed non-collapse versus global independence]\label{rem:typed-non-collapse-vs-independence}
Typed non-collapse is not a global primitive-independence statement. It does not assert that the six roles are logically or
categorically independent in a global algebra; it asserts only that
the pairwise identifications listed as forbidden collapses in
Table~\ref{tab:boundary-matrix} are inadmissible under honest
bookkeeping, leaving the corresponding typed interactions available.
\end{remark}


\subsection{Level-profile architecture}\label{subsec:level-profile-architecture}

Typed non-collapse fixes which identifications are inadmissible; cross-level
claims also require recorded levels and translations. The level trichotomy of
Definition~\ref{def:level-trichotomy} therefore supplies the profile for such
claims.

\begin{definition}[Level profile of an admissible claim]\label{prop:level-profile}
A claim annotation is \emph{cross-level} when records it names are listed in the target fields of level annotations recording different levels; it is a \emph{direct comparison} when
its statement contains an equality or order atom, or an implication, whose two sides read records listed under different recorded levels.
Its \emph{level profile} is a level annotation targeting the claim, together
with, for each such pair of levels, a translation annotation carried in
\(\Audit{D}\) whose source and target contexts are those levels.
Every admissible cross-level claim about a description \(D \in \FATCD{}\) carries a
\emph{level profile}: it locates the claim on the behavioral,
verification, or structural-categorical level, and each cross-level
assertion is accompanied by a translation record in the sense of
Definition~\ref{prop:honest-bookkeeping}. Direct comparison of claims at
different levels without a translation record is inadmissible.
\end{definition}

\subsection{Forgetting and selected lifts}\label{subsec:forgetting-lifts}

Level transitions are controlled by forgetting and selected lifts.

\begin{definition}[Admissible forgetting and selected lifts]\label{prop:forgetting-lifts}
Forgetting from a higher to a lower level (structural-categorical to
verification, verification to behavioral, or structural-categorical directly to
behavioral) is admissible if and only if it carries a record of
its losses carried in \(\Audit{D}\); the loss record may be empty only when
losslessness is proved and recorded. Lifting from a lower level to a higher level is
admissible only when accompanied by
\begin{itemize}[topsep=2pt,itemsep=1pt]
  \item the added data required to support the higher-level statement;
  \item the mark that the lift is selected and non-unique unless uniqueness for its specified source data and target class is proved and recorded;
  \item explicit preservation of lower-level invariants and explicit
    losses where the lift is partial.
\end{itemize}
Lossless forgetting and unique lifts are inadmissible unless the losslessness
or uniqueness is itself proved and recorded.
\end{definition}

\begin{remark}[Lower-to-higher non-determinacy]\label{rem:lower-to-higher-non-determinacy}
Definition~\ref{prop:forgetting-lifts} formalizes the non-determinacy of
lower-level data: given behavioral or verification content alone, the admissibility regime
forces no choice of structural-categorical content, and any asserted choice
must be marked as a selected lift. Where the lower-level data do determine the
higher-level datum, as the descended map of a descending update is determined,
that determinacy is admissible with its recorded proof.
\end{remark}

\subsection{Activation thresholds}\label{subsec:activation-thresholds}

The regime must also state when a projected role is active. The
attachment protocol of Definition~\ref{def:threshold-attachment} is
lifted to a regime-level principle.

\begin{remark}[Activation thresholds]\label{prop:activation-thresholds}
At the regime level, the threshold-attachment discipline of
Definition~\ref{def:threshold-attachment} reads as follows. An
admissible active derived role projection \(\pi_i(D)\) exists only when
\(D\) carries the threshold, witness, level, collapse, nonclaim and audit data
required by Definition~\ref{def:threshold-attachment}. A channel with no
admissible projection receives \(\collapsedbc\) if a failed offer is
recorded, and otherwise \(\belowthr\) if a covered seed is present
(Definition~\ref{def:channel-status}); these are statuses of channels, not of
clusters, and a channel with an admissible projection is \(\activeproj\)
whatever other clusters do. Arbitrary enrichment, adding structure
solely to force a role projection, is inadmissible unless the additions
are themselves honestly bookkept as raw features or annotations. This
restates the attachment protocol at regime level: the sentences before the
enrichment sentence restate Definitions~\ref{def:threshold-attachment}
and~\ref{def:channel-status}, and the enrichment sentence
follows from Definition~\ref{def:fatcd}, since every added record of a member of
\FATCD{} is subject to Definitions~\ref{def:finite}--\ref{def:audited}.
\end{remark}

\subsection{Instrument-relative visibility and suppression}\label{subsec:visibility-suppression}

Threshold-governed activation still leaves open what an instrument makes
visible or suppresses. Admissibility of visibility claims is therefore
instrument-relative.

\begin{definition}[Admissible visibility record]\label{prop:visibility}
Every admissible visibility claim carries an explicit \emph{visibility record}
comprising its visible content, its suppressed content, the instruments
under which the visibility and suppression are asserted, and the
resulting instrument-relative claim strength. Overreading a suppressed
claim as visible under a different instrument is inadmissible. Partial
visibility, in which a claim is visible only up to some coarsening or
translation, is admitted when the visibility record itself makes the
partiality explicit. An \emph{instrument} is a declared record of the
description named in an instrument-visibility annotation's instruments field
or as the measuring instrument of an empirical-bridge claim carried in
\(\Audit{D}\); the \emph{claim strength} is a value of a declared finite ordered
set recorded in that annotation; a visibility claim \emph{asserts
instrument-free truth} when it asserts visible content without naming an
instrument. The \emph{instrument family} of a description is the set of its
instrument records.
\end{definition}

\begin{remark}[No instrument-free truth]\label{rem:no-instrument-free}
Definition~\ref{prop:visibility} does not claim that instruments are
free of distortion or that any single instrument is canonical. It asserts only that admissible visibility claims carry an instrument-relative
visibility record. Instrument completeness is deferred to
Section~\ref{sec:future-work}.
\end{remark}

\subsection{Empirical-bridge admissibility}\label{subsec:empirical-bridge}

Any claim crossing from the closure-algebraic side of the framework
to a substrate-level observable must pass through the
empirical-bridge discipline, related to but narrower than
the empirical-adequacy view of theories~\citep{vanFraassen1980}.

\begin{definition}[Admissible empirical-bridge claim]\label{prop:empirical-bridge}
An empirical-bridge claim about \(D \in \FATCD{}\) is admissible only when
it has, carried in \(\Audit{D}\),
\begin{itemize}[topsep=2pt,itemsep=1pt]
  \item the substrate and observable to which the bridge relates;
  \item the instrument under which the observable is measured;
  \item the level map between the closure-algebraic content and the
    substrate-level content;
  \item a threshold-witness relation recording when the bridge is
    crossed;
  \item the nonclaims suppressed by the bridge record, in particular any
    structural content of \(D\) that is not tracked by the substrate or
    instrument.
\end{itemize}
Empirical realization --- the claim that \(D\) is itself a substrate
observation --- is stronger than empirical-bridge admissibility and is
not asserted by any empirical-bridge claim.
\end{definition}


\subsection{Nonclaim closure across the admissibility regime}\label{subsec:nonclaim-closure}

The rules of this section, from honest bookkeeping
(Definition~\ref{prop:honest-bookkeeping}) to admissible empirical-bridge
claims (Definition~\ref{prop:empirical-bridge}), keep the global
nonclaims of Remark~\ref{rem:atlas-global-nonclaims} in force. The next
proposition lists the overclaims they block.

\begin{proposition}[Admissibility-regime nonclaim closure]\label{prop:nonclaim-closure}
Under the admissibility regime, each of the following claims is
inadmissible; the second kind of claim in the first item cannot occur, since
no member of \FATCD{} carries it:
\begin{itemize}[topsep=2pt,itemsep=1pt]
  \item a claim identifying two of the six roles along a forbidden collapse
    of Table~\ref{tab:boundary-matrix}, or a claim annotation whose statement applies
    a recorded finite map with domain the role names \(P_1,\dots,P_6\) (or
    tuples of them), or an independence predicate on them, whose referenced
    records are only boundary and witness records, and whose statement is not
    represented by their fields in the sense of
    Definition~\ref{def:annotations} (by Definition~\ref{def:audited}(e));
  \item an unconditional empirical-realization claim, that \(D\) is itself a
    substrate observation, where \(D\) records no extension of that
    predicate (criterion after Definition~\ref{def:residual-scheme}), whether
    or not an empirical-bridge record is present;
  \item a lift asserted to be unique without a recorded proof of uniqueness;
  \item a claim that lower-level content determines higher-level content,
    without a recorded proof of that determinacy;
  \item a visibility claim asserting instrument-free truth;
  \item a direct cross-level comparison without a translation record;
  \item a forgetting asserted to be lossless without a recorded proof of
    losslessness;
  \item an empirical-bridge claim lacking any of: a threshold-witness
    relation, a level map, the instrument under which the observable is
    measured, or suppressed-structure nonclaims.
\end{itemize}
\end{proposition}

\begin{proof}
Each requirement of the regime blocks a specific class of overclaims;
we check that it blocks the listed claims.

Honest bookkeeping (Definition~\ref{prop:honest-bookkeeping}) declares
inadmissible every unrecorded claim, implicit promotion, silent
translation, asserted-but-unselected lift, lossless forgetting without
proof, visibility overread, and bare empirical-realization assertion;
these are requirements of Definition~\ref{prop:honest-bookkeeping},
beyond the description-level audit condition of Definition~\ref{def:audited}.

Typed non-collapse (Proposition~\ref{prop:typed-non-collapse}) blocks the
forbidden identifications of Table~\ref{tab:boundary-matrix}. A claim
annotation asserting a full algebra (a map applied to role names) or a global
independence theorem (an independence predicate on role names) must represent
that substantive assertion in fields of its referenced records
(Definition~\ref{def:annotations}). Boundary and \(W_i\)-witness records carry
exactly the fields of their declared checks, and those checks supply neither a
map on role names nor a global-independence certificate. An assertion not represented by the fields of the cited boundary and witness
records requires an additional declared content record. A claim citing only
those records and failing the representation condition cannot be carried by a member of
\FATCD{}. This does not prove that no full
algebra or global independence theorem can exist. Both broader statements remain nonclaims of
Remark~\ref{rem:atlas-global-nonclaims}.

Level-profile architecture (Definition~\ref{prop:level-profile}) blocks
direct cross-level comparison without a translation record, which is
exactly the non-claim that cross-level comparison remains inadmissible
without translation.

Admissible forgetting and selected lifts
(Definition~\ref{prop:forgetting-lifts}) block lossless forgetting,
unique lifts, and lower-to-higher determinacy: forgetting is lossy
unless losslessness is proved and recorded, no lift is admitted as
unique unless uniqueness is proved and recorded, and, by
Remark~\ref{rem:lower-to-higher-non-determinacy}, no structural-categorical
content is asserted to be forced by behavioral or verification content without
a recorded proof. These are
the non-claims that forgetting remains lossy, that no unique-lift
theorem is claimed, and that no lower-to-higher determinacy is claimed.

Activation thresholds (Remark~\ref{prop:activation-thresholds}) block
arbitrary enrichment and threshold-free nontriviality claims, since an
active projection is admissible only with the threshold, witness, level,
collapse, nonclaim and audit data of Definition~\ref{def:threshold-attachment}.

Instrument-relative visibility (Definition~\ref{prop:visibility}) blocks
instrument-free visibility claims and visible-to-suppressed overreads; it does
not impose instrument indexing on every claim kind.

Empirical-bridge admissibility (Definition~\ref{prop:empirical-bridge})
blocks empirical-realization overclaims and any bridge claim lacking a
threshold-witness relation, a level map, the instrument under which the
observable is measured, or suppressed-structure nonclaims. An
empirical-realization claim asserts a predicate, being a substrate
observation, whose extension \(D\) does not record; its truth does not follow
from any finite check of \(D\), so by Definition~\ref{def:audited} no honest
audit certifies it outright, and it is inadmissible whether or not a bridge
record is present. An unrecorded realization predicate may occur as a
hypothesis of an audited conditional consequence; no audit certifies
realization outright or as a consequent. This yields the non-claim that no empirical realization
is claimed and that empirical-bridge claims carry the required data.

Collecting these residues gives each of the listed non-claims. Those that
restate a global nonclaim of Remark~\ref{rem:atlas-global-nonclaims} are
thereby carried forward into the regime; the two further items --- that
forgetting remains lossy absent a recorded proof of losslessness, and that
empirical-bridge claims carry their threshold-witness relation, level map,
instrument and suppressed-structure nonclaims --- are additional disciplines the regime imposes
at this level.
\end{proof}
